

\begin{frame}{Our Approach}

    % 2-agent pursuit-evasion game ... Lambert transfer analysis ... low-thrust analysis ... 
    Given a 2-agent, pursuit-evasion game between satellites, the combined trajectory optimization and game theory problem can be approached with several methods: 
    \begin{itemize}
        \item Intercept, where we only consider the final positions of the satellites 
        \item Rendezvous, where additionally care about the final velocities of the satellites\footnote{For more details on the rendezvous problem, see Paper Appendix.}
    \end{itemize}

\end{frame} 

\begin{frame}{Intercept Geometry}
    For intercept problems, why do we only care about the initial maneuver at $t_1$?

    \begin{columns}
    \begin{column}{0.5\linewidth}
    Lambert's problem provides a solution with two maneuvers: 
    \begin{itemize}
        \item $\boldsymbol{v}_1$ : velocity that Pursuer needs to perform at $t_1$ to arrive at Evader location at $t_2$
        \item $\boldsymbol{v}_2$ : velocity that corrects Pursuer transfer orbit to match Evader velocity at $t_2$
    \end{itemize}
    \end{column}

    \begin{column}{0.5\linewidth}
    \begin{figure}
        \centering
        \includegraphics[width = \textwidth]{pres_images/traj.png}
        \caption{Pursuer and Evader orbits with Lambert solution}
        \label{fig:traj}
    \end{figure}

    \end{column}
    \end{columns}
    
    
\end{frame}

\begin{frame}{Known Simulation Parameters}
    Assume that the following values are given for the pursuer $P$ and the evader $E$: 

    \begin{itemize}
        \item $\boldsymbol{x}^P_0 = [\boldsymbol{r}^P_0, \boldsymbol{v}^P_0]^T$ 
        \item $\boldsymbol{x}^E_0 = [\boldsymbol{r}^E_0, \boldsymbol{v}^E_0]^T$
        \item Dynamics for each player are defined by the Cartesian two-body satellite-Earth system: $\dot{\boldsymbol{x}} = [\boldsymbol{v}, -\frac{\mu \boldsymbol{r}}{|\boldsymbol{r}|^3}]^T$
        \item Earth's gravitational parameter $\mu$
    \end{itemize}

    Time variables: 
    \begin{itemize}
        \item $t_0$: initial simulation time
        \item $t_1$: time at first maneuver
        \item $t_2$: time at second maneuver
        \item Note that $t_0 \neq t_1$, as the initial states of the satellites might not provide the best geometry to perform the maneuver
    \end{itemize}
    
\end{frame}



\begin{frame}{Intercept Problem - Passive Evader}
    We model the problem for the Pursuer as:
    % \begin{align}
    %     \min_{\Delta \boldsymbol{v}_1, t_1, t_2} \;      & |\boldsymbol{r}^P_2 - \boldsymbol{r}^E_2| \\ 
    %     \text{ subject to }             & 
    %                                     & \Delta \boldsymbol{v}_1 = |\boldsymbol{v}^P_1 - \boldsymbol{v}_1| \\
                                        
    % \end{align} 
    % where (P, E) are the pursuer and evader, respectively. 
    % 
    \begin{align}
        \min_{\Delta \boldsymbol{v}_1, t_1, t_2} \;      & |\Delta \boldsymbol{v}_1| \\ 
        \text{ subject to: }             & \boldsymbol{r}^P_2 - \boldsymbol{r}^E_2 = 0 \\
                                        & \dot{\boldsymbol{x}}^P = [\boldsymbol{v}^P, -\frac{\mu \boldsymbol{r}^P}{|\boldsymbol{r}^P|^3}]^T\\
                                        & \dot{\boldsymbol{x}}^E = [\boldsymbol{v}^E, -\frac{\mu \boldsymbol{r}^E}{|\boldsymbol{r}^E|^3}]^T\\
                                        & \Delta \boldsymbol{v}_1 = |\boldsymbol{v}^P_1 - \boldsymbol{v}_1| \\
                                        & \boldsymbol{v}_1, \boldsymbol{v}_2 = Lambert(\boldsymbol{r}^P_0, \boldsymbol{r}^E_0, t_2, t_1, \mu)
    \end{align} 

    $\boldsymbol{r}^E_t$ can be computed by the Pursuer for all times $t$, since the Evader does not perform any maneuvers. This gives us a baseline for what the Pursuer trajectories can look like given the dynamics constraints.
    
\end{frame}

\begin{frame}{Directions to Lambert Solutions}
    Verifying the orbit geometry is necessary in finding solutions using Lambert targeting.

    \begin{columns}
    \begin{column}{0.5\linewidth}
    \begin{figure}
        \centering
        \includegraphics[width = \textwidth]{pres_images/traj.png}
        \caption{Prograde Lambert Solution}
    \end{figure}
    \end{column}

    \begin{column}{0.5\linewidth}
    \begin{figure}
        \centering
        \includegraphics[width = \textwidth]{pres_images/traj_retro.png}
        \caption{Retrograde Lambert Solution}
    \end{figure}
    \end{column}
    \end{columns}

    While the retrograde solution numerically exists, \red{the trajectory \\ goes through Earth!}

\end{frame}

\begin{frame}{Solution Discontinuities - Passive Evader}
    Lambert targeting allows us to compute for numerically feasible $\Delta \boldsymbol{v}$ solutions, but to obtain the minimum, we require iterating through $(t_1, t_2)$ values.

    \begin{columns}
    
        \begin{column}{0.5\linewidth}
        \begin{figure}
            \centering
            \includegraphics[width = 0.9\textwidth]{pres_images/lambert_pro_dvTotalcontour.png}
            \caption{Prograde-Intercept $|\Delta \boldsymbol{v}^P_1|$}
        \end{figure}
        \end{column}
    
        \begin{column}{0.5\linewidth}
        \begin{figure}
            \centering
            \includegraphics[width = 0.8\textwidth]{pres_images/lambert_varyTF_intercept.png}
            \caption{Prograde-Intercept solutions for fixed $t_1$, varying $t_2$}
        \end{figure}
        \end{column}
    
    \end{columns}

    These results led us to constrain the problem to smaller initial orbit differences and a shorter propagation time.
\end{frame}

\begin{frame}{Intercept Problem - Active Evader}
    We reformulate the problem to implicitly find the minimum Lambert solution to focus on minimizing the distance between the satellites when the Evader is allowed to maneuver. 

    The Pursuer's optimization problem now becomes: 
    \begin{align}
        \min_{\Delta \boldsymbol{v}_1, t_1, t_2} \;      & |\boldsymbol{r}^P_2 - \boldsymbol{r}^E_2| \label{eq:min_active_evader_prob} \\ 
        \text{ subject to }             & \dot{\boldsymbol{x}^P} = [\boldsymbol{v}^P, -\frac{\mu \boldsymbol{r}^P}{|\boldsymbol{r}^P|^3}]^T\\
                                        & \dot{\boldsymbol{x}^E} = [\boldsymbol{v}^E, -\frac{\mu \boldsymbol{r}^E}{|\boldsymbol{r}^E|^3}]^T\\
                                        & \Delta \boldsymbol{v}_1 = |\boldsymbol{v}^P_1 - \boldsymbol{v}_1| \\
                                        & \boldsymbol{v}_1, \boldsymbol{v}_2 = minLambert(\boldsymbol{r}^P_0, \boldsymbol{r}^E_0, t_2, t_1, \mu) \label{eq:min_lambert}
    \end{align} 
    
    % \begin{align}
    %     \min_{\Delta \boldsymbol{v}_1, t_1, t_2} \;      & |\Delta \boldsymbol{v}_1| \\ 
    %     \text{ subject to: }             & \boldsymbol{r}^P_2 - \boldsymbol{r}^E_2 = 0 \\
    %                                     & \dot{\boldsymbol{x}}^P = [\boldsymbol{v}^P, -\frac{\mu \boldsymbol{r}^P}{|\boldsymbol{r}^P|^3}]^T\\
    %                                     & \dot{\boldsymbol{x}}^E = [\boldsymbol{v}^E, -\frac{\mu \boldsymbol{r}^E}{|\boldsymbol{r}^E|^3}]^T\\
    %                                     & \Delta \boldsymbol{v}_1 = |\boldsymbol{v}^P_1 - \boldsymbol{v}_1| \\
    %                                     & \boldsymbol{v}_1, \boldsymbol{v}_2 = Lambert(\boldsymbol{r}^P_0, \boldsymbol{r}^E_0, t_2, t_1, \mu)
    % \end{align} 

\end{frame}

% \begin{frame}{Rendezvous Problem} -- PUT IN APPENDIX
%     We model the problem as:
%     \begin{align}
%         \min_{\Delta \boldsymbol{v}_1, \red{\Delta \boldsymbol{v}_2}, t_1, t_2} \;      & |\boldsymbol{r}^P_2 - \boldsymbol{r}^E_2| \\ 
%         \text{ subject to }             & \dot{\boldsymbol{x}^P} = [\boldsymbol{v}^P, -\frac{\mu \boldsymbol{r}^P}{|\boldsymbol{r}^P|^3}]^T\\
%                                         & \dot{\boldsymbol{x}^E} = [\boldsymbol{v}^E, -\frac{\mu \boldsymbol{r}^E}{|\boldsymbol{r}^E|^3}]^T\\
%                                         & \Delta \boldsymbol{v}_1 = |\boldsymbol{v}^P_1 - \boldsymbol{v}_1| \\
%                                         & \red{\Delta \boldsymbol{v}_2 = |\boldsymbol{v}^E_2 - \boldsymbol{v}_2|} \\
%                                         & \boldsymbol{v}_1, \boldsymbol{v}_2 = minLambert(\boldsymbol{r}^P_0, \boldsymbol{r}^E_0, t_2, t_1, \mu)
%     \end{align} 
%     where (P, E) are the pursuer and evader, respectively. 
    
%     % Model the problem as:
%     % \begin{align}
%     %     \min_{\Delta \boldsymbol{v}_1, \red{\Delta \boldsymbol{v}_2}, t_1, t_2} \;      & |\Delta \boldsymbol{v}_1| + \red{|\Delta \boldsymbol{v}_f|}\\ 
%     %     \text{ subject to }             & r^P_2 - r^E_2 = 0.   
%     % \end{align} 

%     % where (P, E) are the pursuer and evader, respectively.

% \end{frame}



\begin{frame}{Short Range Interception} 

    For the short range intercept problem, we assume that we are in an active pursuit game.  
    % and no longer wish to minimize $ \Delta v $ (although it is still subject to constraints). 
% 
    % A reachable set here is described as the set of all possible orbits a given spacecraft can reach given its initial orbit A0 and maximum available $ \Delta V_{max} $. 
% 
    % We define a radius around the evader as the sphere of active pursuit, which we set to be 10 km. 
% 
    The actions for the pursuer are largely unchanged from the longe-range problem, except that we are no longer minimizing $ \Delta v_1 $ (although still subject to constraints), and $ t_1 $ is now fixed at the start.  

    On the other hand, the evader is now actively avoiding interception, and wants to maximize the $ \Delta v $ that the pursuer expends. 
    Inclination changes are the most expensive orbital maneuver satellites can perform, and $ \Delta v $ constraints can be directly translated into $ \Delta i $ constraints. 
    The formula for $ \Delta v $ as a function of desired $ \Delta i $ is  
% 
    \begin{align}
        \Delta i_{mag} & = \sqrt{ v_1^2 + v_2^2 - 2 v_1 v_2 cos( \Delta i ) } \\ 
        \Delta v & = \Delta i_{mag} ( \bar{v}_2 - \bar{v}_1 ) 
    \end{align}

    % v1 = x1[4:6]
    % kep2 = cart2kep(x1, mu)
    % kep2[3] += Δi
    % v2 = kep2cart(kep2, mu)[4:6]
    % Δi_mag = sqrt(norm(v1)^2 + norm(v2)^2 - 2*norm(v1)*norm(v2)*cos(Δi))
    
\end{frame} 

\begin{frame}{Short Range Interception} 

    We set up the simulation as follows: 

    \begin{algorithm}[H]
    \begin{algorithmic}[1]
        \STATE Given initial pursuer state $ \bar{x}_1^P $ and chaser state $ \bar{x}_1^E $ 
        \STATE Pursuer actions: 
        \begin{itemize}
            \item Compute and apply $ \Delta v_1^P $ following Equations \ref{eq:min_active_evader_prob} through \ref{eq:min_lambert} for TOF = 100 s 
            \item Propagate state for 10 s 
        \end{itemize} 
        \STATE Evader actions: 
        \begin{itemize}
            \item Compute and apply $ \Delta v_1^E $ maneuver to maximize inclination change for pursuer 
            \item Propagate state for 10 s 
        \end{itemize} 
        \STATE Repeat steps 2 through 4, setting $ \bar{x}_1^P $ and $ \bar{x}_1^E $ to the final states of the 10 s moving horizon 
        \end{algorithmic}
    \caption{Pursuit Evasion Game}
    \label{alg:seq}
    \end{algorithm}
% 
    % We solve a finite, moving horizon 
    
\end{frame}

\begin{frame}{Pursuit-Evasion Game Animation}

    \centering
    \includemedia[
      activate=pageopen,
      width=240pt,height=150pt,
      addresource=pres_images/pursuit_game.mp4,
      flashvars={%
    src=pres_images/pursuit_game.mp4      % same path as in addresource!
    %scaleMode=stretch % removes black bars
    &autoPlay=true      % optional configuration
    &loop=true          % variables
    } 
    ]{}{StrobeMediaPlayback.swf}

    % \begin{figure}
    %     \centering
    %     \includegraphics{pres_images/pursuit_evasion_game.png}
    %     \caption{Pursuit Evasion Game}
    %     \label{fig:pursuit_evasion}
    % \end{figure}
    
\end{frame}


% \begin{frame}{Long-Range Simulations}

%     Lambert targeting allows us to compute for feasible $\Delta \boldsymbol{v}$ solutions, but to obtain the minimum, we require iterating through $(t_1, t_2)$ values.

%     Plots below show the Pursuer's cost when Evader has:
%     \begin{itemize}
%         \item No control authority, i.e. the Evader stays in the original orbit
%         \item Maneuvers for a limited inclination change (geometric maximum $\Delta \boldsymbol{v}$ due to orbital energy difference
%     \end{itemize}

% \end{frame}
